% ================================================================
% TAILORINIG
% ================================================================
\section{Tailoring}
\label{sec:tailoring}
 
Lo standard ISO/IEC~12207:1995 prevede esplicitamente la possibilità di adattare
(tailoring) il proprio insieme di processi al contesto specifico
del progetto. Il presente documento applica tale adattamento in ragione dei seguenti
fattori contestuali:
 
\begin{itemize}
  \item \textbf{Baseline\textsubscript{G} attuale}: candidatura\textsubscript{G}; nessun
        prodotto software è ancora in sviluppo.
  \item \textbf{Dimensione del gruppo}: sei membri con rotazione documentata dei
        ruoli.
  \item \textbf{Natura didattica}: progetto svolto in ambito accademico con vincoli
        di tempo e risorse predefiniti dal regolamento del corso.
\end{itemize}
 
\subsection{Processi Attivati}
 
La tabella seguente elenca i processi attivati nella fase di candidatura con il
riferimento alla sezione in cui sono normati.
 
\bigskip
\begin{tabularx}{\textwidth}{l l X}
  \toprule
  \textbf{Categoria} & \textbf{Processo} & \textbf{Sezione} \\
  \midrule
  Primario      & Fornitura                     & \ref{sec:fornitura} \\
  Di supporto   & Documentazione                & \ref{sec:documentazione} \\
  Di supporto   & Gestione della Configurazione & \ref{sec:configurazione} \\
  Di supporto   & Verifica                      & \ref{sec:verifica} \\
  Organizzativo & Gestione                      & \ref{sec:gestione} \\
  Organizzativo & Infrastruttura                & \ref{sec:infrastruttura} \\
  \bottomrule
\end{tabularx}
 
\bigskip
I \textbf{processi esclusi} saranno rivalutati a ogni cambio di baseline del progetto e,
ove necessario, introdotti con un aggiornamento del presente documento.
 